% Results section. Every number is a macro or table row generated by paper/make_tables.py
% from artifacts/results_summary.json; nothing here is typed in by hand.

\subsection{The phase boundary on real calibrated views}
Fig.~\ref{fig:phase} and Table~\ref{tab:phase} test Theorem~\ref{thm:main} on WILDTRACK. For each zone redundancy $r$
and colluding set size $k$, we sample real annotated people, let the $k$ Byzantine cameras report a fabricated
position chosen to minimize honest contradiction, and count how often at least one honest view contradicts it. On the
ground plane the empirical boundary sits exactly on $r=k{+}1$ for detection and $r=2k{+}1$ for masking, in every cell
(agreement \AgreeGroundDet{}). Below the boundary the detection rate is zero, as the indistinguishable-worlds
construction requires, and above it an honest witness contradicts the fabrication with rate \WitnessGround{}.

Off the plane the agreement drops to \AgreeOffDet{}. All the disagreeing cells lie one witness \emph{above} the
predicted boundary, never below it. The cause is crowding. When the adversary slides a fabricated head along an honest
camera's ray, a real head belonging to someone else sometimes lies near the fabricated point in another honest view,
so a single non-degenerate witness contradicts the fabrication with rate only \WitnessOff{}. Proposition~\ref{prop:prob}
predicts exactly this shift. With $p=\WitnessOff{}$ and $\delta=0.05$, the required redundancy for $k=1,2,3$ is the
last column of Table~\ref{tab:phase}, one more than the exact bound, and that is where the empirical boundary lies.

\begin{figure*}[t]
\centering
\includegraphics[width=\textwidth]{fig2_byzantine_quorum_bound.png}
\caption{Compression-by-Contradiction phase boundary. (a) Analytical masking redundancy $m(n,k)=2k{+}1$, shown only
where it is feasible ($m\le n$). (b) Empirical detection rate on WILDTRACK ground-plane targets, with the theoretical
detection (solid) and masking (dashed) boundaries. (c) The same for off-plane (head) targets, where the adversary may
exploit epipolar degeneracy along one honest ray.}
\label{fig:phase}
\end{figure*}

\begin{table}[t]
\centering
\caption{Theory versus empirical phase boundary on real calibrated views. ``Agree'' is the fraction of $(r,k)$ cells
whose empirical detection (masking) outcome matches Theorem~\ref{thm:main}; $p$ is the honest-witness contradiction
rate; the last column is Proposition~\ref{prop:prob}'s redundancy for $k=1,2,3$ at $\delta=0.05$.}
\label{tab:phase}
\footnotesize
\setlength{\tabcolsep}{2.5pt}
\begin{tabular}{lccccc}
\toprule
Setting & Cells & \makecell{Agree\\(det.)} & \makecell{Agree\\(mask)} & $p$ & \makecell{$m_{\delta}$\\$k{=}1,2,3$} \\
\midrule
\inputtable{generated/table_phase.tex}
\bottomrule
\end{tabular}
\end{table}

\subsection{Main benchmark}
Table~\ref{tab:main} reports the test episodes (frames 200--399) averaged over five seeds, with the standard deviation
across seeds. The SCSPV attestor has the fewest coverage-weighted blind minutes (\OursBlind{}\,$\pm$\,\OursBlindStd{}
per hour against \BestBaseBlind{} for naive majority vote, a \BlindReductionPct{}\% reduction) and the fewest deceived
minutes (\OursDeceived{} against \MajDeceived{}). The blind-time margin over majority vote is small. The larger
difference is in attestation quality. The attestor's frame AUROC is \OursAUROC{}, against \SingleAUROC{} for
single-camera analytics, \MajAUROC{} for majority vote and \StaticAUROC{} for static redundancy. It also detects the
largest fraction of attacks (\OursDetRate{}) and detects them soonest (\OursDetLat{}\,s after onset on average).

False alarms are where the attestor does not win. Static redundancy raises only \StaticFA{} false alarms per hour,
because it never re-allocates and so rarely acts on a false positive, but it pays for this with \StaticBlind{} blind
minutes per hour. The attestor averages \OursFA{} false alarms per hour with a large spread across seeds
($\pm$\,\OursFAStd{}). On every seed the SCSPV loss selected the highest threshold in the validation grid
($\theta=\OursThetaRange{}$). Fig.~\ref{fig:pareto}(a) shows why. Over the upper half of the grid the attestor's blind
time is flat while its disruption cost keeps falling, so the loss keeps preferring higher thresholds, and the chosen
operating point sits at the edge of the grid, where small differences between seeds move the false-alarm count a lot.
A finer grid near $\theta=1$ is the obvious next step. For thresholds of 0.3 and above, the attestor's frontier lies
below those of single-camera analytics and majority vote. At the two lowest thresholds it quarantines honest cameras
aggressively and its blind time rises above both. No threshold brings it down to static redundancy's false-alarm
rate.

\begin{table*}[t]
\centering
\caption{Main results on the test episodes (mean\,$\pm$\,std over five seeds; best in bold). Blind and deceived time
are coverage-weighted minutes per hour; detection latency is measured from attack onset.}
\label{tab:main}
\footnotesize
\setlength{\tabcolsep}{4pt}
\begin{tabular}{lcccccc}
\toprule
Method & \makecell{Blind\\(min/h)\,$\downarrow$} & \makecell{Deceived\\(min/h)\,$\downarrow$} &
\makecell{False alarms\\(/h)\,$\downarrow$} & \makecell{Detection\\rate\,$\uparrow$} &
\makecell{Detection\\latency (s)\,$\downarrow$} & \makecell{Frame\\AUROC\,$\uparrow$} \\
\midrule
\inputtable{generated/table_main.tex}
\bottomrule
\end{tabular}
\end{table*}

\begin{figure*}[t]
\centering
\includegraphics[width=\textwidth]{fig3_blind_time_pareto.png}
\caption{Availability results over five seeds. (a) Pareto frontier of coverage-weighted blind minutes against
false-alarm disruption cost, traced by sweeping the decision threshold (thin lines: individual seeds). (b) Distribution
of blind minutes per test episode (white dot: mean; bar: interquartile range). (c) Blind time against the compromised
fraction $k/n$ (bands: $\pm$\,one standard deviation across seeds).}
\label{fig:pareto}
\end{figure*}

\subsection{Per-attack and per-$k/n$ behaviour}
Table~\ref{tab:attack} splits the results by attack family. The attestor has the lowest blind time under
deauthentication, RTSP saturation and replay. Replay is where cross-view geometry matters most. A replayed stream has
normal transport telemetry and plausible motion, so single-camera analytics detects it poorly, while the attestor
detects it in \OursDetRateReplay{} of episodes because the replayed detections disagree with live overlapping views.
Static redundancy has by far the highest blind time under replay, plausibly because a replayed witness inside a fixed
two-camera set leaves a tie that majority voting cannot break.

Metadata fabrication is the attestor's weakest case, with a detection rate of only \OursDetRateFab{}, below
single-camera analytics. Every method does poorly here. The adversary starts each phantom in the zone with the fewest
honest witnesses, possibly none, so for part of the episode no honest camera can contradict it. In that regime
Theorem~\ref{thm:main} says that no attestor can detect the fabrication. The bound therefore explains why all four
methods fail here, but it does not explain why the attestor trails single-camera analytics, and we leave that gap
open. Majority vote matches or slightly beats the attestor's blind time on fabrication and mixed attacks. Table~\ref{tab:kn} and Fig.~\ref{fig:pareto}(c) show the dependence on $k/n$. The attestor is best in six
of the seven configurations. Blind time is not monotone in $k/n$, because it depends on how much redundancy the
particular camera subset leaves, which is the quantity the theory identifies, rather than on the ratio alone. The two
peaks are the five-camera subsets, where the analytics budget $B=\lceil 0.4n'\rceil=2$ leaves little slack for hand-off.

\begin{table*}[t]
\centering
\caption{Results by attack family (mean over five seeds). Left: coverage-weighted blind minutes per hour
(mean\,$\pm$\,std, best in bold). Right: fraction of attack episodes detected.}
\label{tab:attack}
\footnotesize
\setlength{\tabcolsep}{4pt}
\begin{tabular}{l cccc cccc}
\toprule
& \multicolumn{4}{c}{Blind (min/h)\,$\downarrow$} & \multicolumn{4}{c}{Detection rate\,$\uparrow$} \\
\cmidrule(lr){2-5}\cmidrule(lr){6-9}
Attack & Single & Majority & Static & Ours & Single & Majority & Static & Ours \\
\midrule
\inputtable{generated/table_attack.tex}
\bottomrule
\end{tabular}
\end{table*}

\begin{table}[t]
\centering
\caption{Blind minutes per hour by compromised fraction $k/n'$ (mean\,$\pm$\,std over five seeds; best in bold).}
\label{tab:kn}
\footnotesize
\setlength{\tabcolsep}{3pt}
\begin{tabular}{lcccc}
\toprule
$k/n'$ & Single & Majority & Static & Ours \\
\midrule
\inputtable{generated/table_kn.tex}
\bottomrule
\end{tabular}
\end{table}

\subsection{Ablations and the untrusted prior}
Table~\ref{tab:ablation} and Fig.~\ref{fig:ablation}(a,b) remove one component at a time. Removing the epipolar and
homography residuals costs the most detection quality (AUROC \AblFullAUROC{} to \AblNoEpiAUROC{}), which confirms that
cross-view geometry carries the attestation signal. Removing transport telemetry costs less (\AblNoTelAUROC{}) but
makes the false-alarm rate higher and less stable. The hand-off policy matters most for availability. Replacing the MILP
with greedy hand-off raises blind time from \AblFullBlind{} to \AblGreedyBlind{} minutes per hour, and disabling
hand-off raises it to \AblNoHandBlind{}.

Removing the ADS-B/Remote-ID prior leaves AUROC unchanged (\AblNoAdsbAUROC{}). The capture was a single overnight
window with sparse traffic, so the prior had little to contribute, and we do not claim a benefit from it. What the
experiment does establish is the safety property. Under the transplanted-traffic spoof, the largest change the prior
made to any camera's risk probability was \SpoofShift{} (per-seed maximum, averaged over seeds). By construction the
prior can never move a logit by more than $2\beta=\SpoofBound{}$, and in practice it stayed far below that.

\begin{table}[t]
\centering
\caption{Ablation (mean\,$\pm$\,std over five seeds). The three hand-off rows reuse the full model's scores, so their
AUROC and false-alarm rate are identical by construction.}
\label{tab:ablation}
\footnotesize
\setlength{\tabcolsep}{2pt}
\resizebox{\columnwidth}{!}{%
\begin{tabular}{lcccc}
\toprule
Variant & AUROC & Blind & Deceived & FA (/h) \\
\midrule
\inputtable{generated/table_ablation.tex}
\bottomrule
\end{tabular}}
\end{table}

\begin{figure*}[t]
\centering
\includegraphics[width=\textwidth]{fig5_ablation_latency.png}
\caption{(a) Frame AUROC and (b) blind minutes per hour when removing one component (error bars: std over five seeds).
(c) Per-frame latency distributions on one CPU thread for the full seven-camera network; the dotted line marks 15\,ms.}
\label{fig:ablation}
\end{figure*}

\subsection{Latency}
Fig.~\ref{fig:ablation}(c) measures each stage on real frames for all seven cameras on a single CPU thread. The
ST-GNN step takes \LatGnnPfifty{}\,ms at the median and \LatGnnPninetyfive{}\,ms at the 95th percentile. With the
epipolar and homography checks, the attestor as a whole takes \LatAttestPfifty{}\,ms (p50) and
\LatAttestPninetyfive{}\,ms (p95), inside a 15\,ms per-frame budget. Two stages fall outside that budget and we report
them separately. An uncached MILP solve takes \LatMilpPfifty{}\,ms (p50) and \LatMilpPninetyfive{}\,ms (p95), but it
runs only when a quarantine decision changes, and repeated coverage signatures are served from a cache. Horn--Schunck
optical flow takes \LatFlowPfifty{}\,ms per camera at the median (\LatFlowPninetyfive{}\,ms at p95) in our unoptimized
CPU implementation. It runs once per camera and so parallelizes across edge devices, and it fits WILDTRACK's 2\,fps
stream, but meeting a 15\,ms budget at full frame rate would need a GPU or a lighter flow estimator.

\subsection{Airport case study: ATL Concourse F}
WILDTRACK is a plaza, so we also place \CaseCams{} cameras over real OpenStreetMap geometry of ATL Concourse F
(\CaseCells{} zones weighted by aeroway class, terminal highest) with an analytics budget of $B=\CaseBudget{}$. The
camera layout is our own, because real layouts are not public. The resulting visibility graph has $\Delta=\CaseDelta{}$,
so Corollary~\ref{cor:degree} caps ground-plane masking at $k\le\CaseKmax{}$. We compromise $k=\CaseK{}$ cameras and
replay the detection delays measured in the benchmark. Fig.~\ref{fig:case} shows the result. Without hand-off the
compromised cameras' zones stay blind for \CaseNoneBlind{} weighted minutes. Greedy hand-off reduces this to
\CaseGreedyBlind{}, and the MILP to \CaseMilpBlind{}. On this layout the two policies are almost identical, because the
budget is large enough that greedy choices are rarely wrong. The ADS-B prior was active over \CaseAdsbCells{} zones and
shifted the normalized zone weights by \CaseSpoofLone{} in $L_1$ under a spoof, within the $[1,1+\beta]$ envelope.

\begin{figure*}[t]
\centering
\includegraphics[width=0.92\textwidth]{fig4_zone_criticality_heatmap.png}
\caption{Case study on real OpenStreetMap geometry of ATL Concourse F (camera layout hypothetical). (a) MILP stream
allocation before the attack; zone colour is criticality $w_z$, outlined squares are unwatched zones. (b) After
$k=\CaseK{}$ cameras are detected and quarantined, the MILP hands their zones to healthy overlapping cameras (green
rings). (c) Coverage-weighted blind weight over time for no hand-off, greedy hand-off and MILP hand-off.}
\label{fig:case}
\end{figure*}
